% ============================================================
% Plantilla: Presentación con Beamer (16:9)
% ============================================================
\documentclass[aspectratio=169]{beamer}

\usepackage[utf8]{inputenc}
\usepackage[T1]{fontenc}
\usepackage[spanish]{babel}
\usepackage{amsmath}
\usepackage{booktabs}

\usetheme{Madrid}
\usecolortheme{seahorse}
\setbeamertemplate{navigation symbols}{}

\title[Título corto]{Título de la presentación}
\subtitle{Subtítulo opcional}
\author{Nombre del Expositor}
\institute{Universidad}
\date{\today}

\begin{document}

\begin{frame}
  \titlepage
\end{frame}

\begin{frame}{Agenda}
  \tableofcontents
\end{frame}

\section{Introducción}
\begin{frame}{Contexto del problema}
  \begin{itemize}
    \item Primer punto importante.
    \item Segundo punto importante.
    \pause
    \item Este punto aparece al avanzar la diapositiva.
  \end{itemize}
\end{frame}

\section{Propuesta}
\begin{frame}{Propuesta en dos columnas}
  \begin{columns}[T]
    \begin{column}{0.48\textwidth}
      \begin{block}{Ventajas}
        \begin{itemize}
          \item Más rápida.
          \item Más segura.
        \end{itemize}
      \end{block}
    \end{column}
    \begin{column}{0.48\textwidth}
      \begin{alertblock}{Limitaciones}
        \begin{itemize}
          \item Requiere conexión.
        \end{itemize}
      \end{alertblock}
    \end{column}
  \end{columns}
\end{frame}

\begin{frame}{Modelo matemático}
  \begin{equation*}
    \hat{y} = \beta_0 + \sum_{i=1}^{n} \beta_i x_i
  \end{equation*}
  \begin{exampleblock}{Interpretación}
    Cada coeficiente $\beta_i$ mide el efecto de la variable $x_i$.
  \end{exampleblock}
\end{frame}

\section{Resultados}
\begin{frame}{Resultados}
  \centering
  \begin{tabular}{lcc}
    \toprule
    Método & Precisión & Tiempo \\
    \midrule
    Anterior & 81\,\% & 2.4 s \\
    Propuesto & \textbf{89\,\%} & \textbf{1.1 s} \\
    \bottomrule
  \end{tabular}
\end{frame}

\begin{frame}
  \centering
  {\Huge ¡Gracias!}\par
  \vspace{1em}
  ¿Preguntas?
\end{frame}

\end{document}
